% short version
Creative thinking is a fundamental aspect of human cognition, and divergent thinking—the capacity to generate novel and varied ideas—is widely regarded as its core generative engine. Large language models (LLMs) have recently demonstrated impressive performance on divergent thinking tests and prior work has shown that models with higher task performance tend to be more aligned to human brain activity. However, existing brain-LLM alignment studies have focused on passive, non-creative tasks. Here, we explore brain alignment during creative thinking using fMRI data from 170 participants performing the Alternate Uses Task (AUT). We extract representations from LLMs varying in size (270M–72B) and measure alignment to brain responses via Representational Similarity Analysis (RSA), targeting the creativity-related default mode and frontoparietal networks. We find that brain-LLM alignment scales with model size (default mode network only) and idea originality (both networks), with effects strongest early in the creative process. We further show that post-training objectives shape alignment in functionally selective ways: a creativity-optimized \texttt{Llama-3.1-8B-Instruct} preserves alignment with high-creativity neural responses while reducing alignment with low-creativity ones; a human behavior fine-tuned model elevates alignment with both; and a reasoning-trained variant shows the opposite pattern, suggesting chain-of-thought training steers representations away from creative neural geometry toward analytical processing. These results demonstrate that post-training objectives selectively reshape LLM representations relative to the neural geometry of human creative thought.

% long version
% Creative thinking is a fundamental aspect of human cognition, and divergent thinking—the capacity to generate novel and varied ideas—is widely regarded as its core generative engine. Large language models (LLMs) have recently demonstrated impressive performance on divergent thinking tests, and prior work has shown that models with higher task performance tend to be more aligned to human brain activity. However, existing brain-LLM alignment studies have focused on passive, non-creative tasks. Here, we explore brain alignment during creative thinking, exploring whether larger or more creative LLMs are more aligned with neural activity recorded during the Alternate Uses Task (AUT)—a standard divergent thinking task requiring the generation of original object uses. We investigate this question using fMRI data from 170 human participants as they generated ideas. We extract representations from a range of LLMs varying in size (from 270M to 72B) and measure their alignment to brain responses using Representational Similarity Analysis (RSA), targeting canonical creativity-related brain regions within the default mode and frontoparietal networks. 
% We find that LLM-brain alignment during creative thinking scales positively with both model size (for the default mode network only) and the rated originality of ideas (for both brain networks). Critically, this pattern is strongest early in the creative process, indicating stage-dependent shifts in alignment. 
% We further find that post-training objectives shape brain alignment in functionally selective ways. A \texttt{Llama-3.1-8B-Instruct} model optimized for creativity via preference optimization preserves moderate alignment with high-creativity neural responses while significantly decreasing alignment with low-creativity ones, whereas a model fine-tuned to simulate human behavior shows elevated alignment with both populations. 
% On the other hand, the same model trained with reasoning chains distilled from \texttt{DeepSeek-R1} exhibits the opposite pattern, suggesting that chain-of-thought reasoning training steers model representations away from the neural geometry of creative ideation and toward more analytical processing. These results demonstrate that post-training objectives can selectively reshape LLM representations relative to the neural geometry of human creative thought in interpretable and distinct ways.

% We also find that post-training a \texttt{Llama-3.1-8B-Instruct} model via preference optimization using creative human responses selectively increases its alignment with brain responses during the generation of highly original ideas, while decreasing brain alignment recorded during less original ideas. These results suggest that post-training for creativity can selectively reshape LLM representations toward the neural geometry of human creative thought. 